\documentclass[conference]{IEEEtran}

% ---------------------------------------------------------------
% Packages
% ---------------------------------------------------------------
\usepackage{cite}
\usepackage{amsmath,amssymb}
\usepackage{booktabs}
\usepackage{graphicx}
\usepackage{xcolor}
\usepackage{url}
\usepackage{microtype}
\usepackage{array}
\usepackage{balance}

% ---------------------------------------------------------------
% Document
% ---------------------------------------------------------------
\begin{document}

\title{Automated Resilience Assessment for\\Autonomous System SoC Architectures}

% Double-blind submission
\author{\IEEEauthorblockN{[Author names omitted for double-blind review]}}

\maketitle

% ---------------------------------------------------------------
\begin{abstract}
Hardware consolidation onto System-on-Chip (SoC) platforms gives autonomous
systems size, weight, and power advantages but creates shared-bus attack
surfaces where a single compromise can propagate to safety-critical
functions. We present an automated resilience assessment that takes an SoC
architectural model with pre-computed security feature assignments and
access-control policies, then systematically generates fault and compromise
scenarios, computes firewall-aware blast radii, detects single points of
failure (SPOFs), and quantifies risk amplification factors aligned with
NIST SP~800-30. The analysis is formulated in Answer Set Programming (ASP)
and runs in under 30 seconds for architectures up to 11 components.
We evaluate on two case studies: an 8-component SoC drone (18 scenarios)
and an 11-component DARPA CASE UAV (24 scenarios). The tool reveals that
topology-driven weaknesses---policy server concentration, common-mode bus
failures, and single-string service chains---persist across all
optimization strategies, demonstrating that these are architectural
vulnerabilities rather than policy-dependent artefacts.
\end{abstract}

\begin{IEEEkeywords}
Resilience Assessment, System-on-Chip, Autonomous Systems, Answer Set
Programming, Zero-Trust Architecture, Fault Analysis
\end{IEEEkeywords}

% ---------------------------------------------------------------
\section{Introduction}

Autonomous systems increasingly rely on SoC platforms that consolidate
sensing, computation, communication, and actuation onto shared bus
interconnects. This consolidation creates cascading failure paths: a
compromised IP core or failed bus can deny service to multiple
safety-critical functions simultaneously. Standards such as
DO-326A~\cite{do326a} require systematic security risk assessment, yet
current practice relies on manual analysis that does not scale with
architectural complexity.

We address this gap with an automated resilience assessment phase that
operates on an SoC graph model. The assessment is the third stage of a
larger design space exploration (DSE) workflow: Phase~1 optimally assigns
security features under FPGA resource constraints; Phase~2 synthesizes
zero-trust access-control policies with firewall placement. Phase~3---the
focus of this paper---consumes both outputs and automatically generates
scenarios covering component compromise, bus failure, control-plane attack,
redundancy-group degradation, and combined fault-plus-compromise events.

For each scenario, the ASP solver computes: (i)~structural and
firewall-aware blast radius, (ii)~service availability and capability
impact, (iii)~risk amplification relative to the baseline, and
(iv)~attack paths up to five hops. The key finding across both case
studies is that the worst-case vulnerabilities are topology-driven---they
persist identically across max-security, min-resource, and balanced
optimization strategies---meaning they cannot be mitigated by feature
selection alone and require architectural redesign.

% ---------------------------------------------------------------
\section{Background and Workflow Context}

\textbf{SoC Graph Model.}
The SoC is modeled as a directed graph $G=(V,E)$ with typed vertices:
bus masters~$V_M$, receivers~$V_R$, bus interconnects~$V_B$, firewalls
(PEPs)~$V_{FW}$, and policy servers (PDPs)~$V_{PS}$. Each component
carries trust-domain, CIA-impact, and exploitability annotations.

\textbf{Phase~1} assigns one security feature and one logging feature per
component, minimizing weighted risk under LUT, flip-flop, and power
constraints with solver-verified optimality~\cite{gebser_asp}. The risk
model follows an additive formulation aligned with NIST
SP~800-30~\cite{nist_sp80030}, avoiding multiplicative artefacts on
ordinal scales.

\textbf{Phase~2} synthesizes least-privilege access-control policies with
firewall placement following NIST SP~800-207 zero-trust
principles~\cite{nist_zta}. It identifies excess privileges and trust
anchor gaps (missing hardware root of trust, attestation, secure boot).

Phase~3 takes the complete annotated graph---with feature assignments,
firewall topology, and ACL policies---as input for resilience analysis.

% ---------------------------------------------------------------
\section{Automated Resilience Assessment}

\subsection{Scenario Generation}

Scenarios are generated systematically from the topology, not manually
enumerated. The generator produces six classes:

\begin{enumerate}
\item \textbf{Single-component compromise}: each master and receiver.
\item \textbf{Bus failure}: each bus interconnect fails, severing all
  connected components.
\item \textbf{Control-plane attack}: policy server compromise (bypasses
  governed firewalls) and PEP bypass.
\item \textbf{Redundancy-group degradation}: progressive member failures
  within quorum groups.
\item \textbf{Combined fault+compromise}: compound events pairing a
  component compromise with a bus failure.
\item \textbf{Control-plane degradation}: all policy servers fail
  simultaneously.
\end{enumerate}

\subsection{Blast Radius and Risk Amplification}

For each scenario~$s$, the solver computes the set of components whose
risk increases due to lost firewall protection, compromised policy
decisions, or severed bus connectivity. The \emph{risk amplification
factor} is:
\begin{equation}
\alpha(s) = \frac{\sum_{c} R_s(c)}{\sum_{c} R_{\mathit{base}}(c)}
\label{eq:amplification}
\end{equation}
where $R_s(c)$ is the per-component risk under scenario~$s$ and
$R_{\mathit{base}}(c)$ is the baseline risk from Phase~1. Amplification
$\alpha > 1$ indicates risk increase; values exceeding $2\times$ flag
critical architectural weaknesses.

\subsection{SPOF and Capability Analysis}

The solver maps each scenario to affected \emph{services} (groups of
components providing a function) and \emph{capabilities} (mission-level
functions such as flight control or navigation). A scenario that disables
$k \ge 3$ capabilities simultaneously is flagged as a critical SPOF.
Attack paths of up to five hops are enumerated from compromised
components to safety-critical targets.

% ---------------------------------------------------------------
\section{Case Study Results}

We evaluate on two architectures targeting the Xilinx Zynq-7000
(PYNQ-Z2) FPGA platform, each analyzed under multiple Phase~1
optimization strategies.

\subsection{SoC Drone (TC9)}

The TC9 architecture has 8 components, 2 bus masters, 2 NoC buses, a
5-member redundancy group, and 1 safety-critical component. Phase~1
produces solver-verified optimal assignments under three strategies
(max-security: 74 weighted risk at 17\% LUTs; min-resources: 101 risk at
10\% LUTs). Phase~2 places 2 firewalls and 1 policy server, identifying
9 excess privileges and 6 trust anchor gaps.

Phase~3 generates 18 scenarios. The two critical findings are:
(1)~\texttt{ps0} compromise yields $2.5\times$ risk amplification because
this single policy server governs both firewalls---a control-plane
concentration SPOF; (2)~\texttt{noc0} failure is a common-mode event that
simultaneously disables all 5 redundancy-group members, defeating the
availability benefit of redundancy. Both findings are invariant across
all three optimization strategies.

\subsection{DARPA CASE UAV}

The UAV architecture, translated from the DARPA CASE surveillance
system~\cite{darpa_case,hasan_saecps}, has 11 components, 3 masters,
4 buses, 4 safety-critical components, and zero redundancy groups.
Phase~1 achieves risk~=~0 under max-security (34\% LUTs) and
risk~=~39 under min-resources (15\% LUTs), both solver-verified optimal.
Phase~2 places 1 firewall and 1 policy server, flagging 22 excess
privileges and 10 trust anchor gaps.

Phase~3 generates 24 scenarios. Under min-resources (baseline risk~42.2),
\texttt{ps\_uart} compromise produces $2.06\times$ amplification
(risk~87.0). The most severe availability finding is
\texttt{bus\_mc\_failure}, which loses 5 capabilities simultaneously
(flight control, navigation, OTA update, attestation, ground comms)---a
catastrophic single-string SPOF that no security feature assignment can
address. Under max-security, all 24 scenario risks are zero (base risk~0
$\times$ amplification = 0), yet the topology-driven SPOF persists:
\texttt{bus\_mc} failure still loses the same 5 capabilities.

\begin{table}[t]
\caption{Phase~3 Resilience Summary: Both Case Studies}
\label{tab:resilience}
\centering
\scriptsize
\setlength{\tabcolsep}{3pt}
\begin{tabular}{lcc}
\toprule
Metric & TC9 (8-comp) & UAV (11-comp) \\
\midrule
Scenarios generated          & 18          & 24 \\
Worst scenario               & ps0 comp.   & ps\_uart comp. \\
Worst amplification          & $2.50\times$& $2.06\times$ \\
Max capabilities lost (1 event) & 5 (noc0) & 5 (bus\_mc) \\
Control-plane SPOFs          & 1 (ps0)     & 1 (ps\_uart) \\
Common-mode bus failures     & 1 (noc0)    & 1 (bus\_mc) \\
Invariant across strategies  & Yes         & Yes \\
Solve time                   & $<$15\,s    & $<$30\,s \\
\bottomrule
\end{tabular}
\end{table}

\subsection{Cross-Study Observations}

Table~\ref{tab:resilience} summarizes both case studies. Despite different
sizes, topologies, and redundancy characteristics, both architectures
exhibit the same structural pattern: a single policy server whose
compromise amplifies risk above $2\times$, and a single bus whose failure
disables the majority of mission capabilities. These findings are
\emph{strategy-invariant}---they appear identically under max-security,
min-resources, and balanced optimization---confirming they are topological
properties of the architecture rather than artefacts of any particular
security policy configuration.

The TC9 redundancy group illustrates a subtler failure mode: while
quorum-based redundancy improves availability under independent component
failures, the shared \texttt{noc0} bus creates a common-mode path that
defeats the redundancy assumption entirely.

% ---------------------------------------------------------------
\section{Related Work}

Hardware isolation on SoCs through TrustZone and secure interconnects
provides static partitioning but not automated vulnerability
assessment~\cite{fiorin_noc}. The DARPA CASE program~\cite{darpa_case}
produced AADL-based model transformations with CakeML-verified
components; Hasan et~al.~\cite{hasan_saecps} applied manual ZTA patterns
to the same UAV system but did not perform systematic resilience
analysis. Zero-trust principles for avionics were advocated
in~\cite{ztas_dasc}, and NIST SP~800-207~\cite{nist_zta} provides the
policy architecture framework. Our contribution is automated,
solver-based scenario generation and blast-radius computation that
complements these manual approaches.

% ---------------------------------------------------------------
\section{Conclusion}

We presented an automated resilience assessment for SoC-based autonomous
systems that generates fault and compromise scenarios from the
architectural topology, computes firewall-aware blast radii, and detects
single points of failure. Evaluation on an 8-component drone and an
11-component DARPA CASE UAV shows that the worst-case
vulnerabilities---policy server concentration and single-string bus
dependencies---are topology-driven and persist across all optimization
strategies. This demonstrates that security feature optimization alone is
insufficient; architectural redundancy at the bus and control-plane
levels is required for resilient autonomous system SoC designs. Future
work includes multi-chip topology extension, runtime-adaptive mode
transitions, and integration with DO-326A certification evidence
generation.

% ---------------------------------------------------------------
\balance
\bibliographystyle{IEEEtran}
\bibliography{rasacc_references}

\end{document}
